\documentclass[a4paper]{arrowhead}

\usepackage[yyyymmdd]{datetime}
\usepackage{etoolbox}
\usepackage[utf8]{inputenc}
\usepackage[T1]{fontenc}
\usepackage{multirow}
\usepackage{longtable}
\usepackage{tikz}
\usetikzlibrary{positioning,shapes.multipart,arrows.meta}

\renewcommand{\dateseparator}{-}
\setlength{\emergencystretch}{3em}
\lstset{language=HTTP,breakatwhitespace=false,columns=fullflexible}

\newcommand{\fparam}[1]{\textit{\textcolor{ArrowheadBlue}{#1}}}

%% Special references
\newcommand{\mref}[2][]{{\textcolor{ArrowheadPurple}{\hyperref[sec:model:#2]{\ifstrempty{#1}{#2}{#1}}}}}
\newcommand{\oref}[1]{{\textcolor{ArrowheadBlue}{\hyperref[sec:operations:#1]{#1}}}}
\newcommand{\pdef}[1]{{\textcolor{ArrowheadGrey}{#1 \label{sec:model:primitives:#1} \label{sec:model:primitives:#1s}}}}
\newcommand{\pref}[1]{{\textcolor{ArrowheadGrey}{\hyperref[sec:model:primitives:#1]{#1}}}}

% as template_idd.tex; the input and output arguments are passed already formatted (\mref or \pref)
\newrobustcmd\osubsection[5]{
  \addtocounter{subsection}{1}
  \addcontentsline{toc}{subsection}{\protect\numberline{\thesubsection}operation \textcolor{ArrowheadBlue}{#1}}
  \subsection*{
    \thesubsection\quad
    #2 \textcolor{ArrowheadPurple}{#3} \\
    \small
    \hspace*{0.075\textwidth}\begin{minipage}{0.1\textwidth}
      \vspace*{1mm}
      Operation: \\
      \notblank{#4}{Input: \\}{}
      \notblank{#5}{Output: \\}{}
    \end{minipage}
    \begin{minipage}{0.825\textwidth}
      \vspace*{1mm}
      \textcolor{ArrowheadBlue}{#1} \\
      \notblank{#4}{#4 \\}{}
      \notblank{#5}{#5 \\}{}
    \end{minipage}
  }
  \label{sec:operations:#1}
}
\newrobustcmd\msubsection[2]{
  \addtocounter{subsection}{1}
  \addcontentsline{toc}{subsection}{\protect\numberline{\thesubsection}#1 \textcolor{ArrowheadPurple}{#2}}
  \subsection*{\thesubsection\quad#1 \textcolor{ArrowheadPurple}{#2}}
  \label{sec:model:#2} \label{sec:model:#2s}
}
%%

\begin{document}
\sloppy

%% Arrowhead Document Properties
\ArrowheadTitle{certificateAuthorityManagement HTTP/TLS/JSON}
\ArrowheadServiceID{certificateAuthorityManagement}
\ArrowheadType{Interface Design Description}
\ArrowheadTypeShort{IDD}
\ArrowheadVersion{5.2.2}
\ArrowheadDate{2026-10-30}
\ArrowheadAuthor{Jerker Delsing}
\ArrowheadStatus{RELEASE CANDIDATE}
\ArrowheadContact{jerker.delsing@ltu.se}
\ArrowheadFooter{\href{www.arrowhead.eu}{www.arrowhead.eu}}
\ArrowheadSetup
%%

%% Front Page
\begin{center}
  \vspace*{1cm}
  \huge{\arrowtitle}

  \vspace*{0.2cm}
  \LARGE{\arrowtype}
  \vspace*{1cm}

  \Large{Service ID: \textit{"\arrowid"}}

  \vspace*{\fill}

  % Front Page Image

  \vspace*{1cm}
  \vspace*{\fill}

  % Front Page Abstract
  \begin{abstract}
    This document describes the HTTP/1.1 protocol with TLS payload security and JSON payload encoding variant of the \texttt{certificate\allowbreak{}Authority\allowbreak{}Management} service of the CertificateAuthority system.
  \end{abstract}

  \vspace*{1cm}
\end{center}

\newpage
%%

%% Table of Contents
\tableofcontents
\newpage
%%

\section{Overview}
\label{sec:overview}

This document describes the \texttt{certificate\allowbreak{}Authority\allowbreak{}Management} service interface of the CertificateAuthority system, which provides operator management of the certificate register and of trusted keys, including revocation.

It's implemented using protocol, encoding as stated in below table:
\begin{table}[ht!]
  \centering
  \begin{tabular}{|l|l|l|}
    \hline
    \rowcolor{gray!33} Profile type & Type & Version \\ \hline
    Transfer protocol & HTTP & 1.1 \cite{rfc9112} \\ \hline
    Data encryption & TLS & 1.3 \cite{rfc8446}, 1.2 \cite{rfc5246} \\ \hline
    Encoding & JSON & RFC 8259 \cite{rfc8259} \\ \hline
    Compression & N/A & - \\ \hline
    Semantics & N/A & - \\ \hline
    Ontology & N/A & - \\ \hline
  \end{tabular}
  \caption{Communication and semantics details used for the \texttt{certificate\allowbreak{}Authority\allowbreak{}Management} service interface}
  \label{tab:comunication_semantics_profile}
\end{table}

This document provides the Interface Design Description IDD to the \textit{certificateAuthorityManagement -- Service Description} document \cite{cau-certificateAuthorityManagement-sd}.
For further details about how this service is meant to be used, please consult that document.

The rest of this document describes how to realize the \texttt{certificate\allowbreak{}Authority\allowbreak{}Management} service
interface in detail. Both in terms of its operations (Section \ref{sec:operations}) and its information model (Section \ref{sec:model}).

\subsection{Interface registration}

The service is registered in the ServiceRegistry with the interface template \texttt{generic\_\allowbreak{}https}, policy \texttt{CERT\_\allowbreak{}AUTH}, access address \texttt{certificate-\allowbreak{}authority}, port 8448 and base path \texttt{/\allowbreak{}certificate-\allowbreak{}authority}. The \texttt{operations} property of the interface lists the operations below with their method and path relative to the base path (path parameters are placeholders such as \texttt{\{id\}} that the consumer fills in):

\noindent\begin{longtable}{| p{5.2cm} | p{1.6cm} | p{8.4cm} |} \hline
\rowcolor{gray!33} Operation & Method & Path \\ \hline \endhead
\texttt{list-\allowbreak{}certificates} & GET & \texttt{/\allowbreak{}mgmt/\allowbreak{}certificates} \\ \hline
\texttt{revoke-\allowbreak{}certificate} & DELETE & \texttt{/\allowbreak{}mgmt/\allowbreak{}certificates/\allowbreak{}\{id\}} \\ \hline
\texttt{revoke-\allowbreak{}certificate-\allowbreak{}by-\allowbreak{}content} & POST & \texttt{/\allowbreak{}mgmt/\allowbreak{}revocations} \\ \hline
\texttt{list-\allowbreak{}keys} & GET & \texttt{/\allowbreak{}mgmt/\allowbreak{}keys} \\ \hline
\texttt{add-\allowbreak{}key} & PUT & \texttt{/\allowbreak{}mgmt/\allowbreak{}keys} \\ \hline
\texttt{delete-\allowbreak{}key} & DELETE & \texttt{/\allowbreak{}mgmt/\allowbreak{}keys/\allowbreak{}\{id\}} \\ \hline
\end{longtable}

\newpage

\section{Service Operations}
\label{sec:operations}

The interfaces of the \texttt{certificate\allowbreak{}Authority\allowbreak{}Management} service, its operations, data models and implementation are provided below.
In particular, each subsection first names the HTTP method and path used to call the operation, after which it names the abstract operation from the \texttt{certificate\allowbreak{}Authority\allowbreak{}Management} SD document, as well as input and output types.

Every request needs mutual TLS with a system or operator certificate of the cloud and \texttt{Authorization: Bearer IDENTITY-TOKEN//<token>}; the token is verified with Authentication and must belong to the system named by the certificate. Request and response bodies are \texttt{application/json} unless stated otherwise.
Errors are answered with \texttt{\{\"error\": \"<text>\"\}} and the status codes \texttt{400 Bad Request} (invalid input), \texttt{401 Unauthorized} (no valid identity), \texttt{403 Forbidden} (caller not permitted), \texttt{404 Not Found}, \texttt{500 Internal Server Error}, \texttt{502 Bad Gateway} (a system it depends on failed) and \texttt{503 Service Unavailable} (overloaded).

\osubsection{list-certificates}{GET}{/\allowbreak{}certificate-authority/\allowbreak{}mgmt/\allowbreak{}certificates}{}{\pref{Array}$<$\allowbreak{}\mref{CertificateRecord}$>$}

Returns the certificate register, newest first.

A successful call returns \texttt{200 OK} with the output in the body.

\begin{lstlisting}[label={lst:list-certificates},caption={An \oref{list-certificates} invocation.}]
GET /certificate-authority/mgmt/certificates HTTP/1.1
Host: certificate-authority:8448
Authorization: Bearer IDENTITY-TOKEN//<token>
\end{lstlisting}

\begin{lstlisting}[label={lst:list-certificates-resp},caption={An \oref{list-certificates} response.}]
[
  {
    "id": 42,
    "serial_number": "4189327123456789",
    "common_name": "Sensor1.TestCloud.Company.arrowhead.eu",
    "created_by": "Onboarding",
    "valid_from": "2026-10-22T12:00:00Z",
    "valid_until": "2027-10-22T12:00:00Z",
    "certificate_pem": "-----BEGIN CERTIFICATE-----\\n..."
  }
]
\end{lstlisting}

\osubsection{revoke-certificate}{DELETE}{/\allowbreak{}certificate-authority/\allowbreak{}mgmt/\allowbreak{}certificates/\allowbreak{}\fparam{\{id\}}}{\pref{String}}{\mref{RemovalResult}}

Revokes a register record by identifier (\texttt{\{id\}}) and republishes the CRL.

Path parameter: \fparam{id}.

A successful call returns \texttt{200 OK} with the output in the body.

\begin{lstlisting}[label={lst:revoke-certificate},caption={An \oref{revoke-certificate} invocation.}]
DELETE /certificate-authority/mgmt/certificates/<id> HTTP/1.1
Host: certificate-authority:8448
Authorization: Bearer IDENTITY-TOKEN//<token>
\end{lstlisting}

\begin{lstlisting}[label={lst:revoke-certificate-resp},caption={An \oref{revoke-certificate} response.}]
{
  "success": true
}
\end{lstlisting}

\osubsection{revoke-certificate-by-content}{POST}{/\allowbreak{}certificate-authority/\allowbreak{}mgmt/\allowbreak{}revocations}{\mref{RevocationRequest}}{\mref{RevocationResult}}

Revokes any certificate of this CA given by content; unknown ones are added to the register as revoked.

A successful call returns \texttt{200 OK} with the output in the body.

\begin{lstlisting}[label={lst:revoke-certificate-by-content},caption={An \oref{revoke-certificate-by-content} invocation.}]
POST /certificate-authority/mgmt/revocations HTTP/1.1
Host: certificate-authority:8448
Authorization: Bearer IDENTITY-TOKEN//<token>
Content-Type: application/json

{
  "certificate": "<PEM or Base64 DER>"
}
\end{lstlisting}

\begin{lstlisting}[label={lst:revoke-certificate-by-content-resp},caption={An \oref{revoke-certificate-by-content} response.}]
{
  "id": 42,
  "serialNumber": "4189327123456789",
  "revoked": true,
  "changed": true
}
\end{lstlisting}

\osubsection{list-keys}{GET}{/\allowbreak{}certificate-authority/\allowbreak{}mgmt/\allowbreak{}keys}{}{\pref{Array}$<$\allowbreak{}\mref{TrustedKeyRecord}$>$}

Returns the trusted keys.

A successful call returns \texttt{200 OK} with the output in the body.

\begin{lstlisting}[label={lst:list-keys},caption={An \oref{list-keys} invocation.}]
GET /certificate-authority/mgmt/keys HTTP/1.1
Host: certificate-authority:8448
Authorization: Bearer IDENTITY-TOKEN//<token>
\end{lstlisting}

\begin{lstlisting}[label={lst:list-keys-resp},caption={An \oref{list-keys} response.}]
[
  {
    "id": 42,
    "public_key_pem": "-----BEGIN PUBLIC KEY-----\\n...",
    "description": "Display may read",
    "valid_after": "string",
    "valid_before": "string",
    "created_at": "string"
  }
]
\end{lstlisting}

\osubsection{add-key}{PUT}{/\allowbreak{}certificate-authority/\allowbreak{}mgmt/\allowbreak{}keys}{\mref{TrustedKeyRequest}}{\mref{TrustedKeyResult}}

Adds a trusted key with validity.

A successful call returns \texttt{201 Created} with the output in the body.

\begin{lstlisting}[label={lst:add-key},caption={An \oref{add-key} invocation.}]
PUT /certificate-authority/mgmt/keys HTTP/1.1
Host: certificate-authority:8448
Authorization: Bearer IDENTITY-TOKEN//<token>
Content-Type: application/json

{
  "publicKey": "-----BEGIN PUBLIC KEY-----\\n...",
  "description": "Display may read",
  "validAfter": "2026-10-22T12:00:00Z",
  "validBefore": "2026-10-22T12:00:00Z"
}
\end{lstlisting}

\begin{lstlisting}[label={lst:add-key-resp},caption={An \oref{add-key} response.}]
{
  "id": 42,
  "trusted": true
}
\end{lstlisting}

\osubsection{delete-key}{DELETE}{/\allowbreak{}certificate-authority/\allowbreak{}mgmt/\allowbreak{}keys/\allowbreak{}\fparam{\{id\}}}{\pref{String}}{\mref{RemovalResult}}

Deletes a trusted key by identifier (\texttt{\{id\}}).

Path parameter: \fparam{id}.

A successful call returns \texttt{200 OK} with the output in the body.

\begin{lstlisting}[label={lst:delete-key},caption={An \oref{delete-key} invocation.}]
DELETE /certificate-authority/mgmt/keys/<id> HTTP/1.1
Host: certificate-authority:8448
Authorization: Bearer IDENTITY-TOKEN//<token>
\end{lstlisting}

\begin{lstlisting}[label={lst:delete-key-resp},caption={An \oref{delete-key} response.}]
{
  "success": true
}
\end{lstlisting}

\newpage

\section{Data Models}
\label{sec:model}

Here, all data objects that can be part of the service calls associated with this service are listed in alphabetic order.
Note that each subsection, which describes one type of object, begins with the \textit{struct} keyword, which is meant to denote a JSON \pref{Object} that must contain certain fields, or names, with values conforming to explicitly named types.
Fields that are not needed may be omitted or \texttt{null} in requests; empty fields are omitted in responses.
As a complement to the primary types defined in this section, there is also a list of secondary types in Section \ref{sec:model:primitives}, which are used to represent things like hashes, identifiers and texts.

\msubsection{struct}{CertificateRecord}




\noindent\begin{tabularx}{\textwidth}{| >{\raggedright\arraybackslash}p{3.9cm} | >{\raggedright\arraybackslash\small}p{5.3cm} | X |} \hline
\rowcolor{gray!33} Object Field & Value Type & Description \\ \hline
"id" & \pref{Number} & Record identifier. \\ \hline
"serial\_number" & \pref{String} & Serial number. \\ \hline
"common\_name" & \pref{String} & Subject common name. \\ \hline
"created\_by" & \pref{String} & Requesting system (\texttt{revoked-\allowbreak{}by:\allowbreak{}\textless{}caller\textgreater{}} for imported revocations). \\ \hline
"valid\_from" & \pref{String} & Start of validity. \\ \hline
"valid\_until" & \pref{String} & End of validity. \\ \hline
"revoked\_at" & \pref{String} & Revocation time, null if not revoked. \\ \hline
"certificate\_pem" & \pref{String} & The certificate, PEM. \\ \hline
\end{tabularx}



\msubsection{struct}{RemovalResult}




\noindent\begin{tabularx}{\textwidth}{| >{\raggedright\arraybackslash}p{3.9cm} | >{\raggedright\arraybackslash\small}p{5.3cm} | X |} \hline
\rowcolor{gray!33} Object Field & Value Type & Description \\ \hline
"success" & \pref{Boolean} & Whether a record was changed; 404 when false. \\ \hline
\end{tabularx}



\msubsection{struct}{RevocationRequest}




\noindent\begin{tabularx}{\textwidth}{| >{\raggedright\arraybackslash}p{3.9cm} | >{\raggedright\arraybackslash\small}p{5.3cm} | X |} \hline
\rowcolor{gray!33} Object Field & Value Type & Description \\ \hline
"certificate" & \pref{String} & Certificate to revoke, PEM or Base64 DER; must be issued by this cloud CA. \\ \hline
\end{tabularx}



\msubsection{struct}{RevocationResult}




\noindent\begin{tabularx}{\textwidth}{| >{\raggedright\arraybackslash}p{3.9cm} | >{\raggedright\arraybackslash\small}p{5.3cm} | X |} \hline
\rowcolor{gray!33} Object Field & Value Type & Description \\ \hline
"id" & \pref{Number} & Register record of the certificate. \\ \hline
"serial\allowbreak{}Number" & \pref{String} & Serial number. \\ \hline
"revoked" & \pref{Boolean} & Always true. \\ \hline
"changed" & \pref{Boolean} & False if it was already revoked. \\ \hline
\end{tabularx}



\msubsection{struct}{TrustedKeyRecord}




\noindent\begin{tabularx}{\textwidth}{| >{\raggedright\arraybackslash}p{3.9cm} | >{\raggedright\arraybackslash\small}p{5.3cm} | X |} \hline
\rowcolor{gray!33} Object Field & Value Type & Description \\ \hline
"id" & \pref{Number} & Record identifier. \\ \hline
"public\_key\_pem" & \pref{String} & The key, PEM. \\ \hline
"description" & \pref{String} & Description. \\ \hline
"valid\_after" & \pref{String} & Start of validity. \\ \hline
"valid\_before" & \pref{String} & End of validity. \\ \hline
"created\_at" & \pref{String} & Creation time. \\ \hline
\end{tabularx}



\msubsection{struct}{TrustedKeyRequest}




\noindent\begin{tabularx}{\textwidth}{| >{\raggedright\arraybackslash}p{3.9cm} | >{\raggedright\arraybackslash\small}p{5.3cm} | X |} \hline
\rowcolor{gray!33} Object Field & Value Type & Description \\ \hline
"public\allowbreak{}Key" & \pref{String} & PEM Subject\allowbreak{}Public\allowbreak{}Key\allowbreak{}Info (alias \texttt{publicKeyPem}). \\ \hline
"description" & \pref{String} & Description. \\ \hline
"valid\allowbreak{}After" & \pref{DateTime} & Start of validity; default now. \\ \hline
"valid\allowbreak{}Before" & \pref{DateTime} & End of validity; default now plus \texttt{CERT\_\allowbreak{}VALIDITY\_\allowbreak{}DAYS}. \\ \hline
\end{tabularx}



\msubsection{struct}{TrustedKeyResult}




\noindent\begin{tabularx}{\textwidth}{| >{\raggedright\arraybackslash}p{3.9cm} | >{\raggedright\arraybackslash\small}p{5.3cm} | X |} \hline
\rowcolor{gray!33} Object Field & Value Type & Description \\ \hline
"id" & \pref{Number} & Record identifier. \\ \hline
"trusted" & \pref{Boolean} & Always true. \\ \hline
\end{tabularx}



\subsection{Primitives}
\label{sec:model:primitives}

As all messages are encoded using the JSON format \cite{rfc8259}, the following primitive constructs, part of that standard, become available.
Note that the official standard is defined in terms of parsing rules, while this list only concerns syntactic information.
Furthermore, the \pref{Object} and \pref{Array} types are given optional generic type parameters, which are used in this document to signify when pair values or elements are expected to conform to certain types.

\noindent\begin{tabularx}{\textwidth}{| p{4.6cm} | X |} \hline
\rowcolor{gray!33} Type & Description \\ \hline
\pdef{Array} & A JSON Array \cite{rfc8259}; \pref{Array}$<$A$>$ has elements of type A. \\ \hline
\pdef{Boolean} & JSON Boolean, \texttt{true} or \texttt{false}. \\ \hline
\pdef{DateTime} & JSON String with an ISO 8601 UTC timestamp, e.g. \texttt{2026-10-22T12:00:00Z}. \\ \hline
\pdef{Number} & JSON Number, an integer. \\ \hline
\pdef{Object} & JSON Object; \pref{Object}$<$A$>$ has string keys and values of type A. \\ \hline
\pdef{String} & JSON String. \\ \hline
\end{tabularx}

\newpage

\bibliographystyle{IEEEtran}
\bibliography{bibliography}

\newpage

\section{Revision History}
\subsection{Amendments}

\noindent\begin{tabularx}{\textwidth}{| p{1cm} | p{3cm} | p{2cm} | X | p{4cm} |} \hline
\rowcolor{gray!33} No. & Date & Version & Subject of Amendments & Author \\ \hline

1 & 2026-10-22 & \arrowversion & Document created from the Arrowhead LaTeX templates for the 5.2.2 Java Spring implementation, including the local extensions of this workspace. & Jerker Delsing \\ \hline
2 & 2026-10-30 & \arrowversion & Checked against the primary Java source tree (Eclipse-Arrowhead-java): services, operations, paths and data models unchanged; missing primitive definitions added. & Claude Opus 5.5 (draft) \\ \hline
\end{tabularx}

\subsection{Quality Assurance}

\noindent\begin{tabularx}{\textwidth}{| p{1cm} | p{3cm} | p{2cm} | X |} \hline
\rowcolor{gray!33} No. & Date & Version & Approved by \\ \hline

1 & & \arrowversion  &  \\ \hline

\end{tabularx}

\end{document}
